\begin{lemma}
Let $\Omega\subset\mathbb{C}$ be open, let $H\colon\mathbb{C}\to\mathbb{C}$ be
$C^\infty$ on $\Omega$, let $s,k\ge0$ be integers, and let
$P\in\mathbb{P}(s,k)$ relative to $H$ in the sense of \noderef{PClass}.  Then
there exist functions $Q_1,Q_2\in\mathbb{P}(s+1,k)$ relative to $H$ such that
\[
  \frac{\partial P}{\partial\xi}(\xi)=Q_1(\xi)
  \quad\text{and}\quad
  \frac{\partial P}{\partial\bar\xi}(\xi)=Q_2(\xi)
  \qquad\text{for every }\xi\in\Omega,
\]
with $\partial/\partial\xi$ as in \noderef{DZ} and $\partial/\partial\bar\xi$
as in \noderef{DZbar}.
\end{lemma}
\begin{proof}
By \noderef{PClass} fix an integer $r\ge0$, complex numbers $c_1,\dots,c_r$ and
lists $M_1,\dots,M_r$ of derivative symbols with
\begin{equation}\tag{1}
  s(M_\alpha)=s,\qquad k(M_\alpha)=k\qquad(\alpha=1,\dots,r),
\end{equation}
\begin{equation}\tag{2}
  P(\xi)=\sum_{\alpha=1}^{r}c_\alpha H^{M_\alpha}(\xi)\qquad\text{for every }\xi\in\mathbb{C}.
\end{equation}

\emph{Step 1: each symbol value is $C^\infty$ on $\Omega$.}  Let
$d=(i,j,\beta)$ be a derivative symbol, and put
\[
  g_d=\frac{\partial^{i+j}H}{\partial\xi^i\partial\bar\xi^j},
\]
the iterated derivative of \noderef{WirtingerIter}.  By part (a) of
\noderef{WirtingerIterSwap}, $g_d$ is $C^\infty$ on $\Omega$.  By
\noderef{SymbolValue}, $H^d$ is either $g_d$ or $\overline{g_d}$; complex
conjugation is an $\mathbb{R}$-linear isometry of $\mathbb{C}$, so composing a
$C^\infty$ function on $\Omega$ with it again gives a $C^\infty$ function on
$\Omega$.  Hence $H^d$ is $C^\infty$ on $\Omega$, and in particular $g_d$ and
$H^d$ are real-differentiable at every point of $\Omega$.

\emph{Step 2: differentiating one symbol value raises its order by one.}  Let
$d=(i,j,\beta)$ be a derivative symbol and define two triples
\[
  d'=\begin{cases}(i+1,\,j,\,\beta)&\text{if $\beta$ is false},\\
                  (i,\,j+1,\,\beta)&\text{if $\beta$ is true},\end{cases}
  \qquad
  d''=\begin{cases}(i,\,j+1,\,\beta)&\text{if $\beta$ is false},\\
                   (i+1,\,j,\,\beta)&\text{if $\beta$ is true}.\end{cases}
\]
In each of the four cases the two indices are non-negative with positive sum,
so $d'$ and $d''$ are derivative symbols in the sense of
\noderef{DerivSymbol}, and
\begin{equation}\tag{3}
  |d'|=|d|+1,\qquad |d''|=|d|+1 .
\end{equation}
We claim that for every $\xi\in\Omega$
\begin{equation}\tag{4}
  \frac{\partial H^{d}}{\partial\xi}(\xi)=H^{d'}(\xi),
  \qquad
  \frac{\partial H^{d}}{\partial\bar\xi}(\xi)=H^{d''}(\xi).
\end{equation}

Suppose first that $\beta$ is false, so $H^d=g_d$ by \noderef{SymbolValue}.
By \noderef{WirtingerIter}, $g_d=(\partial/\partial\xi)^i(\partial/\partial\bar\xi)^jH$,
so applying $\partial/\partial\xi$ once gives
$(\partial/\partial\xi)^{i+1}(\partial/\partial\bar\xi)^jH$, which is
$\partial^{i+1+j}H/\partial\xi^{i+1}\partial\bar\xi^{j}$ and hence, by
\noderef{SymbolValue}, equals $H^{(i+1,j,\text{false})}=H^{d'}$; this is an
identity of functions on all of $\mathbb{C}$.  For the second identity, part
(b) of \noderef{WirtingerIterSwap} gives, for $\xi\in\Omega$,
\[
  \frac{\partial g_d}{\partial\bar\xi}(\xi)
  =\frac{\partial^{i+j+1}H}{\partial\xi^{i}\partial\bar\xi^{j+1}}(\xi)
  =H^{(i,j+1,\text{false})}(\xi)=H^{d''}(\xi),
\]
the middle equality by \noderef{SymbolValue}.

Suppose now that $\beta$ is true, so $H^d=\overline{g_d}$ by
\noderef{SymbolValue}.  By Step 1, $g_d$ is real-differentiable at every
$\xi\in\Omega$, so \noderef{WirtingerConjRule} applies to $g_d$ at each such
$\xi$ and gives
\[
  \frac{\partial\overline{g_d}}{\partial\xi}(\xi)
  =\overline{\frac{\partial g_d}{\partial\bar\xi}(\xi)},
  \qquad
  \frac{\partial\overline{g_d}}{\partial\bar\xi}(\xi)
  =\overline{\frac{\partial g_d}{\partial\xi}(\xi)} .
\]
For the first, part (b) of \noderef{WirtingerIterSwap} gives
$\dfrac{\partial g_d}{\partial\bar\xi}(\xi)
 =\dfrac{\partial^{i+j+1}H}{\partial\xi^{i}\partial\bar\xi^{j+1}}(\xi)$, whence
by \noderef{SymbolValue}
\[
  \frac{\partial H^{d}}{\partial\xi}(\xi)
  =\overline{\frac{\partial^{i+j+1}H}{\partial\xi^{i}\partial\bar\xi^{j+1}}(\xi)}
  =H^{(i,j+1,\text{true})}(\xi)=H^{d'}(\xi).
\]
For the second, $\partial g_d/\partial\xi
=\partial^{i+1+j}H/\partial\xi^{i+1}\partial\bar\xi^{j}$ on all of $\mathbb{C}$
by \noderef{WirtingerIter}, as in the case $\beta$ false, whence by
\noderef{SymbolValue}
\[
  \frac{\partial H^{d}}{\partial\bar\xi}(\xi)
  =\overline{\frac{\partial^{i+1+j}H}{\partial\xi^{i+1}\partial\bar\xi^{j}}(\xi)}
  =H^{(i+1,j,\text{true})}(\xi)=H^{d''}(\xi).
\]
This proves (4) in all cases.

\emph{Step 3: Leibniz on a monomial.}  Let $M=(d_1,\dots,d_k)$ be a list of
derivative symbols.  For $1\le t\le k$ let
\[
  M[t]=(d_1,\dots,d_{t-1},\,d_t',\,d_{t+1},\dots,d_k),
  \qquad
  M\langle t\rangle=(d_1,\dots,d_{t-1},\,d_t'',\,d_{t+1},\dots,d_k)
\]
be the lists obtained from $M$ by replacing its $t$-th entry $d_t$ by $d_t'$,
respectively by $d_t''$, in the notation of Step 2.  Then
\begin{equation}\tag{5}
  k(M[t])=k(M),\qquad k(M\langle t\rangle)=k(M),
  \qquad
  s(M[t])=s(M\langle t\rangle)=s(M)+1 .
\end{equation}
Indeed, by \noderef{MonomialValue} $k(\,\cdot\,)$ is the length of the list,
and replacing one entry does not change the length; and $s(\,\cdot\,)$ is the
sum of the orders of the entries, of which exactly one summand changes, from
$|d_t|$ to $|d_t'|=|d_t|+1$ respectively $|d_t''|=|d_t|+1$, by (3).

We claim that for every $\xi\in\Omega$
\begin{equation}\tag{6}
  \frac{\partial H^{M}}{\partial\xi}(\xi)=\sum_{t=1}^{k}H^{M[t]}(\xi),
  \qquad
  \frac{\partial H^{M}}{\partial\bar\xi}(\xi)=\sum_{t=1}^{k}H^{M\langle t\rangle}(\xi),
\end{equation}
and prove both by induction on the length $k$ of $M$, the two arguments being
word for word the same: the first uses the first identity of (4), the first
identity of \noderef{WirtingerProductRule} and the operator
$\partial/\partial\xi$ of \noderef{DZ}, and the second uses the second
identity of (4), the second identity of \noderef{WirtingerProductRule} and the
operator $\partial/\partial\bar\xi$ of \noderef{DZbar}.  We write out the
first.

If $k=0$, then $H^{M}$ is the constant function $1$ by
\noderef{MonomialValue}.  A constant function is real-differentiable with zero
derivative at every point, so by \noderef{DZ}
$\partial H^{M}/\partial\xi=0$, which is the empty sum on the right of (6).

Let $k\ge1$ and write $M=(d_1)*M''$ with $M''=(d_2,\dots,d_k)$, where $*$
denotes concatenation.  By \noderef{MonomialValue},
$H^{M}=H^{d_1}\cdot H^{M''}$ pointwise.  By Step 1 every symbol value is
real-differentiable at every point of $\Omega$; hence so is $H^{d_1}$, and so
is the finite pointwise product $H^{M''}$.  Therefore
\noderef{WirtingerProductRule} applies at each $\xi\in\Omega$ and gives
\[
  \frac{\partial H^{M}}{\partial\xi}(\xi)
  =\frac{\partial H^{d_1}}{\partial\xi}(\xi)\,H^{M''}(\xi)
   +H^{d_1}(\xi)\,\frac{\partial H^{M''}}{\partial\xi}(\xi).
\]
By the first identity of (4), the first term equals
$H^{d_1'}(\xi)H^{M''}(\xi)$, which by \noderef{MonomialValue} is
$H^{M[1]}(\xi)$, because $M[1]=(d_1')*M''$ and the monomial of a concatenation
is the product of the two monomials.  By the induction hypothesis applied to
$M''$, whose length is $k-1$,
\[
  H^{d_1}(\xi)\,\frac{\partial H^{M''}}{\partial\xi}(\xi)
  =H^{d_1}(\xi)\sum_{t=2}^{k}H^{\,M''[t-1]}(\xi)
  =\sum_{t=2}^{k}H^{M[t]}(\xi),
\]
again by \noderef{MonomialValue}, since $M[t]=(d_1)*M''[t-1]$ for
$2\le t\le k$.  Adding the two contributions gives the first identity of (6).

\emph{Step 4: assembling $Q_1$ and $Q_2$.}  Define, for $\xi\in\mathbb{C}$,
\[
  Q_1(\xi)=\sum_{\alpha=1}^{r}\sum_{t=1}^{k}c_\alpha\,H^{M_\alpha[t]}(\xi),
  \qquad
  Q_2(\xi)=\sum_{\alpha=1}^{r}\sum_{t=1}^{k}c_\alpha\,H^{M_\alpha\langle t\rangle}(\xi).
\]
By (1) and (5), each of the $rk$ lists $M_\alpha[t]$ and each of the $rk$ lists
$M_\alpha\langle t\rangle$ has $s(\,\cdot\,)=s+1$ and $k(\,\cdot\,)=k$, i.e.\
type $(s+1,k)$.  Hence the coefficients $c_\alpha$ together with these lists
present $Q_1$, respectively $Q_2$, as a finite $\mathbb{C}$-linear combination
of monomials of type $(s+1,k)$ at every point of $\mathbb{C}$, so
$Q_1,Q_2\in\mathbb{P}(s+1,k)$ relative to $H$ by \noderef{PClass}.

Finally let $\xi\in\Omega$.  By (2) the function $P$ agrees at every point of
$\mathbb{C}$ with $\zeta\mapsto\sum_{\alpha}c_\alpha H^{M_\alpha}(\zeta)$, so
the two have the same real Fr\'echet derivative at $\xi$ and hence, by
\noderef{DZ} and \noderef{DZbar}, the same $\partial/\partial\xi$ and
$\partial/\partial\bar\xi$ there.  Each $H^{M_\alpha}$ is real-differentiable
at $\xi$, being a finite product of functions differentiable there by Step 1.
The real Fr\'echet derivative is additive on functions differentiable at $\xi$
and satisfies $D(cf)(\xi)=c\,Df(\xi)$ for a constant $c\in\mathbb{C}$; since
by \noderef{DZ} and \noderef{DZbar} each of $\partial/\partial\xi$ and
$\partial/\partial\bar\xi$ at $\xi$ is a fixed $\mathbb{C}$-linear expression
in the values of that derivative at the two vectors $1$ and $i$, both
operators are likewise additive at $\xi$ and commute with multiplication by
constants.  Therefore, using (6) applied to each $M_\alpha$, which is
legitimate because $k(M_\alpha)=k$ by (1),
\[
  \frac{\partial P}{\partial\xi}(\xi)
  =\sum_{\alpha=1}^{r}c_\alpha\frac{\partial H^{M_\alpha}}{\partial\xi}(\xi)
  =\sum_{\alpha=1}^{r}c_\alpha\sum_{t=1}^{k}H^{M_\alpha[t]}(\xi)
  =Q_1(\xi),
\]
and in the same way
\[
  \frac{\partial P}{\partial\bar\xi}(\xi)
  =\sum_{\alpha=1}^{r}c_\alpha\sum_{t=1}^{k}H^{M_\alpha\langle t\rangle}(\xi)
  =Q_2(\xi).
\]
\end{proof}
